\chapter*{Wstęp}
\addcontentsline{toc}{chapter}{Wstęp}

Współczesne systemy bezzałogowych statków powietrznych (UAV) ulegają dynamicznej transformacji, ewoluując od pojedynczych, zdalnie sterowanych jednostek w kierunku autonomicznych rojów zdolnych do realizacji złożonych misji rozproszonych. Rosnąca liczba agentów operujących we wspólnej przestrzeni powietrznej wymusza porzucenie tradycyjnego modelu sterowania, w którym jeden operator kontroluje jedną maszynę, na rzecz paradygmatu interakcji człowiek-rój (ang. \textit{Human-Swarm Interaction}, HSI). Niestety, klasyczne interfejsy sterowania wprowadzają drastyczne obciążenie poznawcze operatora, co przy zarządzaniu rojem prowadzi do natychmiastowej utraty świadomości sytuacyjnej oraz spadku zaufania do automatyzacji. Gwałtowny rozwój wielkich modeli językowych (LLM) stwarza bezprecedensową okazję do zaprojektowania w pełni naturalnych, głosowych interfejsów komunikacyjnych (HRI), które działając jako semantyczny bufor między człowiekiem a algorytmami rojowymi, pozwalają na wydawanie wysokopoziomowych dyrektyw zamiast mikrozarządzania trajektoriami.

Głównym celem niniejszej pracy jest opracowanie, implementacja i ewaluacja systemu adaptacyjnego sterowania rojem dronów. System ten integruje interfejs głosowy wspierany modelami LLM z modułem dynamicznej selekcji algorytmów nawigacyjnych. Zdefiniowana w pracy hipoteza badawcza zakłada, że włączenie frameworku agentowego (RAI) integrującego LLM ze środowiskiem ROS2 i systemem Nav2, połączone z adaptacyjnymi algorytmami inteligencji stadnej (m.in. DPSO), pozwala na redukcję złożoności sterowania z poziomu liniowego $O(n)$ do poziomu stałego $O(1)$ z punktu widzenia operatora, przy jednoczesnym zwiększeniu efektywności i bezpieczeństwa misji.

Szczegółowy zakres pracy obejmuje zaprojektowanie architektury informatycznej łączącej platformę \texttt{ds-crazyflies} z warstwą percepcyjną i decyzyjną. Moduł decyzyjny będzie autonomicznie oceniał stan otoczenia i na bazie poleceń głosowych dobierał odpowiednią strategię nawigacyjną. 

Struktura pracy została podzielona na cztery zasadnicze rozdziały. Rozdział pierwszy stanowi rygorystyczny przegląd literatury z zakresu barier kognitywnych, interfejsów HRI, koncepcji Embodied AI oraz optymalizacji rojowej. W rozdziale drugim scharakteryzowano zastosowane technologie, w tym ROS2, framework RAI oraz stos nawigacyjny Nav2. Rozdział trzeci, stanowiący inżynierski rdzeń pracy, opisuje autorską architekturę i implementację adaptacyjnego systemu decyzyjnego. Rozdział czwarty zawiera analizę wyników przeprowadzonych badań symulacyjnych w zróżnicowanych scenariuszach operacyjnych, weryfikując przyjęte założenia projektowe.